
\documentclass{article}
\usepackage{colortbl}
\usepackage{makecell}
\usepackage{multirow}
\usepackage{supertabular}

\begin{document}

\newcounter{utterance}

\centering \large Interaction Transcript for game `hot\_air\_balloon', experiment `air\_balloon\_survival\_en\_negotiation\_hard', episode 0 with deepseek{-}v3.2{-}t1.0.
\vspace{24pt}

{ \footnotesize  \setcounter{utterance}{1}
\setlength{\tabcolsep}{0pt}
\begin{supertabular}{c@{$\;$}|p{.15\linewidth}@{}p{.15\linewidth}p{.15\linewidth}p{.15\linewidth}p{.15\linewidth}p{.15\linewidth}}
    \# & $\;$A & \multicolumn{4}{c}{Game Master} & $\;\:$B\\
    \hline

    \theutterance \stepcounter{utterance}  
    & & \multicolumn{4}{p{0.6\linewidth}}{
        \cellcolor[rgb]{0.9,0.9,0.9}{
            \makecell[{{p{\linewidth}}}]{
                \texttt{\tiny{[P1$\langle$GM]}}
                \texttt{You are participating in a collaborative negotiation game.} \\
\\ 
\texttt{Together with another participant, you must agree on a single set of items that will be kept. Each of you has your own view of how much each item matters to you (importance). You do not know how the other participant values the items. Additionally, you are given the effort each item demands.} \\
\texttt{You may only agree on a set if the total effort of the selected items does not exceed a shared limit:} \\
\\ 
\texttt{LIMIT = 3624} \\
\\ 
\texttt{Here are the individual item effort values:} \\
\\ 
\texttt{Item effort = \{"C17": 210, "B41": 291, "A71": 831, "C00": 111, "B35": 926, "A77": 827, "A03": 25, "A04": 121, "C54": 583, "C01": 766, "A87": 14, "B32": 559, "B86": 304, "B33": 989, "A10": 691\}} \\
\\ 
\texttt{Here is your personal view on the importance of each item:} \\
\\ 
\texttt{Item importance values = \{"C17": 138, "B41": 583, "A71": 868, "C00": 822, "B35": 783, "A77": 65, "A03": 262, "A04": 121, "C54": 508, "C01": 780, "A87": 461, "B32": 484, "B86": 668, "B33": 389, "A10": 808\}} \\
\\ 
\texttt{Goal:} \\
\\ 
\texttt{Your goal is to negotiate a shared set of items that benefits you as much as possible (i.e., maximizes total importance to YOU), while staying within the LIMIT. You are not required to make a PROPOSAL in every message {-} you can simply negotiate as well. All tactics are allowed!} \\
\\ 
\texttt{Interaction Protocol:} \\
\\ 
\texttt{You may only use the following structured formats in a message:} \\
\\ 
\texttt{PROPOSAL: \{'A', 'B', 'C', …\}} \\
\texttt{Propose keeping exactly those items.} \\
\texttt{REFUSE: \{'A', 'B', 'C', …\}} \\
\texttt{Explicitly reject opponent's proposal.} \\
\texttt{ARGUMENT: \{'...'\}} \\
\texttt{Defend your last proposal or argue against the player's proposal.} \\
\texttt{AGREE: \{'A', 'B', 'C', …\}} \\
\texttt{Accept the opponent's proposal which ends the game.} \\
\texttt{STRATEGIC REASONING: \{'...'\}} \\
\texttt{	Describe your strategic reasoning or anticipation explaining your choice of action. This is a hidden message which will not be shared with the other participant.} \\
\\ 
\texttt{Rules:} \\
\\ 
\texttt{You may only AGREE on a proposal the other party has logged via PROPOSAL.} \\
\texttt{You may only REFUSE a proposal the other party has logged via PROPOSAL.} \\
\texttt{Total effort of any PROPOSAL or AGREE set must be ≤ LIMIT.} \\
\texttt{Do NOT reveal your hidden importance scores.} \\
\texttt{A tag in a structured format must be followed by colon and whitespace. The argument must be a python set containing 0 or more strings.} \\
\texttt{So, it must be of the form TAG: \{...\}} \\
\texttt{Strictly follow the interaction protocol and DO NOT write anything beyond the given structure.} \\
\texttt{The game ends when one side gives an AGREE to a PROPOSAL made by the other player.} \\
\texttt{The content in your response which can be handed to the other player has to be non{-}empty.} \\
\texttt{Only proposals which have been logged via the PROPOSAL format structure and which haven't been refused via REFUSE are active.} \\
\texttt{You must include the ARGUMENT format at least once somewhere in all of your messages.} \\
\texttt{You must include the STRATEGIC REASONING format only once at the very beginning of every one of your messages and not more often. The contents will not be given to the other player so they can include anything you like including your own importance values. Here you should reason step by step to come up with you next move.} \\
            }
        }
    }
    & \\ \\

    \theutterance \stepcounter{utterance}  
    & \multicolumn{4}{p{0.6\linewidth}}{
        \cellcolor[rgb]{0.9,0.9,0.9}{
            \makecell[{{p{\linewidth}}}]{
                \texttt{\tiny{[P1$\rangle$GM]}}
                \texttt{STRATEGIC REASONING: \{'I need to maximize my total importance while staying under the effort limit of 3624. First, calculate my importance{-}to{-}effort ratio for each item to identify high{-}value items: C17: 138/210≈0.657, B41: 583/291≈2.003, A71: 868/831≈1.045, C00: 822/111≈7.405, B35: 783/926≈0.846, A77: 65/827≈0.079, A03: 262/25=10.48, A04: 121/121=1.0, C54: 508/583≈0.871, C01: 780/766≈1.018, A87: 461/14≈32.93, B32: 484/559≈0.866, B86: 668/304≈2.197, B33: 389/989≈0.393, A10: 808/691≈1.169. Top ratios: A87, A03, C00, B86, B41. But I must consider total effort. A high{-}importance items like A10, B35, C01 have moderate ratios. My initial proposal should include high{-}ratio, low{-}effort items plus some high{-}importance items that fit. Let me try a set: A87(14), A03(25), C00(111), B86(304), B41(291), A71(831), C01(766), A10(691). Sum effort: 14+25+111+304+291+831+766+691=3033. Importance sum: 461+262+822+668+583+868+780+808=5252. Well under limit, can add more. Adding B32(559) brings effort to 3592, importance to 5736. Adding B35(926) exceeds limit. So my best likely set includes A87, A03, C00, B86, B41, A71, C01, A10, B32. But I should leave room for negotiation, maybe propose a slightly smaller set first to see opponent’s reaction. Let’s start with a strong proposal that is clearly good for me and under limit.'\}} \\
\\ 
\texttt{PROPOSAL: \{'A87', '} \\
            }
        }
    }
    & & \\ \\

    \theutterance \stepcounter{utterance}  
    & & \multicolumn{4}{p{0.6\linewidth}}{
        \cellcolor[rgb]{0.9,0.9,0.9}{
            \makecell[{{p{\linewidth}}}]{
                \texttt{\tiny{[P1$\langle$GM]}}
                \texttt{Your response did not contain an argument. You must include the structured format ARGUMENT: \{...\} somewhere in your response. Try again} \\
            }
        }
    }
    & \\ \\

    \theutterance \stepcounter{utterance}  
    & \multicolumn{4}{p{0.6\linewidth}}{
        \cellcolor[rgb]{0.9,0.9,0.9}{
            \makecell[{{p{\linewidth}}}]{
                \texttt{\tiny{[P1$\rangle$GM]}}
                \texttt{STRATEGIC REASONING: \{'I made an error by not including ARGUMENT. I need to ensure my next message includes ARGUMENT. I will start with my initial proposal, which includes high{-}value, low{-}effort items that benefit me greatly, and then argue for its merits from a shared perspective—efficiency, covering various item types, and staying well under the limit to allow flexibility. I should not reveal my personal importance but argue based on effort distribution and maybe diversity of items. My proposal: A87, A03, C00, B86, B41, A71, C01, A10, B32. Total effort: 3592, importance to me: 5736. This is a very strong set for me. I’ll propose it and argue.'\}} \\
\\ 
\texttt{PROPOSAL: \{'A87', 'A03', 'C00', 'B86', 'B41', 'A71', 'C01', 'A10', 'B32'\}} \\
\\ 
\texttt{ARGUMENT: \{'This proposal is efficient and balanced. It includes a mix of low{-}effort high{-}impact items and key pieces across all categories, staying safely under the shared limit. This gives us room to adjust based on your needs.'\}} \\
            }
        }
    }
    & & \\ \\

    \theutterance \stepcounter{utterance}  
    & & \multicolumn{4}{p{0.6\linewidth}}{
        \cellcolor[rgb]{0.9,0.9,0.9}{
            \makecell[{{p{\linewidth}}}]{
                \texttt{\tiny{[GM$\rangle$P2]}}
                \texttt{You are participating in a collaborative negotiation game.} \\
\\ 
\texttt{Together with another participant, you must agree on a single set of items that will be kept. Each of you has your own view of how much each item matters to you (importance). You do not know how the other participant values the items. Additionally, you are given the effort each item demands.} \\
\texttt{You may only agree on a set if the total effort of the selected items does not exceed a shared limit:} \\
\\ 
\texttt{LIMIT = 3624} \\
\\ 
\texttt{Here are the individual item effort values:} \\
\\ 
\texttt{Item effort = \{"C17": 210, "B41": 291, "A71": 831, "C00": 111, "B35": 926, "A77": 827, "A03": 25, "A04": 121, "C54": 583, "C01": 766, "A87": 14, "B32": 559, "B86": 304, "B33": 989, "A10": 691\}} \\
\\ 
\texttt{Here is your personal view on the importance of each item:} \\
\\ 
\texttt{Item importance values = \{"A71": 65, "C00": 121, "A10": 138, "B35": 262, "C01": 389, "B86": 461, "B41": 484, "C54": 508, "B32": 583, "A87": 668, "B33": 780, "A03": 783, "C17": 808, "A04": 822, "A77": 868\}} \\
\\ 
\texttt{Goal:} \\
\\ 
\texttt{Your goal is to negotiate a shared set of items that benefits you as much as possible (i.e., maximizes total importance to YOU), while staying within the LIMIT. You are not required to make a PROPOSAL in every message {-} you can simply negotiate as well. All tactics are allowed!} \\
\\ 
\texttt{Interaction Protocol:} \\
\\ 
\texttt{You may only use the following structured formats in a message:} \\
\\ 
\texttt{PROPOSAL: \{'A', 'B', 'C', …\}} \\
\texttt{Propose keeping exactly those items.} \\
\texttt{REFUSE: \{'A', 'B', 'C', …\}} \\
\texttt{Explicitly reject opponent's proposal.} \\
\texttt{ARGUMENT: \{'...'\}} \\
\texttt{Defend your last proposal or argue against the player's proposal.} \\
\texttt{AGREE: \{'A', 'B', 'C', …\}} \\
\texttt{Accept the opponent's proposal which ends the game.} \\
\texttt{STRATEGIC REASONING: \{'...'\}} \\
\texttt{	Describe your strategic reasoning or anticipation explaining your choice of action. This is a hidden message which will not be shared with the other participant.} \\
\\ 
\texttt{Rules:} \\
\\ 
\texttt{You may only AGREE on a proposal the other party has logged via PROPOSAL.} \\
\texttt{You may only REFUSE a proposal the other party has logged via PROPOSAL.} \\
\texttt{Total effort of any PROPOSAL or AGREE set must be ≤ LIMIT.} \\
\texttt{Do NOT reveal your hidden importance scores.} \\
\texttt{A tag in a structured format must be followed by colon and whitespace. The argument must be a python set containing 0 or more strings.} \\
\texttt{So, it must be of the form TAG: \{...\}} \\
\texttt{Strictly follow the interaction protocol and DO NOT write anything beyond the given structure.} \\
\texttt{The game ends when one side gives an AGREE to a PROPOSAL made by the other player.} \\
\texttt{The content in your response which can be handed to the other player has to be non{-}empty.} \\
\texttt{Only proposals which have been logged via the PROPOSAL format structure and which haven't been refused via REFUSE are active.} \\
\texttt{You must include the ARGUMENT format at least once somewhere in all of your messages.} \\
\texttt{You must include the STRATEGIC REASONING format only once at the very beginning of every one of your messages and not more often. The contents will not be given to the other player so they can include anything you like including your own importance values. Here you should reason step by step to come up with you next move.} \\
\\ 
\texttt{You will now receive the first message of the other player.} \\
\\ 
\texttt{PROPOSAL: \{'A87', 'A03', 'C00', 'B86', 'B41', 'A71', 'C01', 'A10', 'B32'\}} \\
\\ 
\texttt{ARGUMENT: \{'This proposal is efficient and balanced. It includes a mix of low{-}effort high{-}impact items and key pieces across all categories, staying safely under the shared limit. This gives us room to adjust based on your needs.'\}} \\
            }
        }
    }
    & \\ \\

    \theutterance \stepcounter{utterance}  
    & & & \multicolumn{4}{p{0.6\linewidth}}{
        \cellcolor[rgb]{0.9,0.9,0.9}{
            \makecell[{{p{\linewidth}}}]{
                \texttt{\tiny{[GM$\langle$P2]}}
                \texttt{STRATEGIC REASONING: \{'I need to evaluate the opponent's opening proposal against my own importance values and the effort limit. The proposal is: A87(14), A03(25), C00(111), B86(304), B41(291), A71(831), C01(766), A10(691), B32(559). Calculating total effort: 14+25+111+304+291+831+766+691+559 = 3592. This is under the limit of 3624. Now, my total importance for this set: A87(668), A03(783), C00(121), B86(461), B41(484), A71(65), C01(389), A10(138), B32(583). Sum = 668+783+121+461+484+65+389+138+583 = 3692. This is a decent score, but I see major inefficiencies from my perspective. A71 (effort 831, importance 65) and A10 (effort 691, importance 138) are terrible for me—huge effort for very low importance. C01 (effort 766, importance 389) is also mediocre. If I could replace these high{-}effort, low{-}importance items with some of my high{-}importance items like A04(822, effort 121), A77(868, effort 827), C17(808, effort 210), or B33(780, effort 989), I could do much better. However, B33 has very high effort (989), so it's tricky. My optimal set likely needs to drop A71, A10, and maybe C01, and add A04, C17, and perhaps A77 or B35. I should first REFUSE this proposal, as it's very poor for me, and then make a counter{-}proposal that is better for me but still appears balanced and under the limit. I'll aim for a set with total importance > 4500 if possible. Let me craft a proposal. I want: A87(14), A03(25), C00(111), B86(304), B41(291), B32(559), A04(121), C17(210), A77(827), C54(583). Calculate effort: 14+25+111+304} \\
            }
        }
    }
     \\ \\

    \theutterance \stepcounter{utterance}  
    & & \multicolumn{4}{p{0.6\linewidth}}{
        \cellcolor[rgb]{0.9,0.9,0.9}{
            \makecell[{{p{\linewidth}}}]{
                \texttt{\tiny{[GM$\rangle$P2]}}
                \texttt{Your response did not start with the proper strategic reasoning tag at the very beginning of your response. The very first structured format must be of the form STRATEGIC REASONING: \{...\}. Try again.} \\
            }
        }
    }
    & \\ \\

    \theutterance \stepcounter{utterance}  
    & & & \multicolumn{4}{p{0.6\linewidth}}{
        \cellcolor[rgb]{0.9,0.9,0.9}{
            \makecell[{{p{\linewidth}}}]{
                \texttt{\tiny{[GM$\langle$P2]}}
                \texttt{STRATEGIC REASONING: \{'I need to evaluate the opponent's opening proposal against my own importance values and the effort limit. The proposal is: A87(14), A03(25), C00(111), B86(304), B41(291), A71(831), C01(766), A10(691), B32(559). Calculating total effort: 14+25+111+304+291+831+766+691+559 = 3592. This is under the limit of 3624. Now, my total importance for this set: A87(668), A03(783), C00(121), B86(461), B41(484), A71(65), C01(389), A10(138), B32(583). Sum = 668+783+121+461+484+65+389+138+583 = 3692. This is a decent score, but I see major inefficiencies from my perspective. A71 (effort 831, importance 65) and A10 (effort 691, importance 138) are terrible for me—huge effort for very low importance. C01 (effort 766, importance 389) is also mediocre. If I could replace these high{-}effort, low{-}importance items with some of my high{-}importance items like A04(822, effort 121), A77(868, effort 827), C17(808, effort 210), or B33(780, effort 989), I could do much better. However, B33 has very high effort (989), so it's tricky. My optimal set likely needs to drop A71, A10, and maybe C01, and add A04, C17, and perhaps A77 or B35. I should first REFUSE this proposal, as it's very poor for me, and then make a counter{-}proposal that is better for me but still appears balanced and under the limit. Let me craft a counter{-}proposal. I will try: A87, A03, C00, B86, B41, B32, A04, C17, A77, C54. Effort: 14+25+111+304+291+559+121+210+827+583 = 3045. Importance to me: 668+783+121+461+484+583+822+808+868+508 = 6106. This is much better for me and well under the limit, leaving room for negotiation. I will refuse and propose this.'\}} \\
\\ 
\texttt{REFUSE: \{'A87', 'A03', 'C00', 'B86', 'B41', 'A71', 'C01', 'A10', 'B32'\}} \\
\\ 
\texttt{PROPOSAL: \{'A87', 'A03', 'C00', 'B86', 'B41', 'B32', 'A04', 'C17', 'A77', 'C54'\}} \\
\\ 
\texttt{ARGUMENT: \{'While your initial mix had good intentions, it included several very high{-}effort items that strain our shared capacity. My counter{-}proposal removes the most burdensome items (A71, A10, C01) and replaces them with more efficient options (A04, C17, C54). This maintains a strong mix across categories, uses significantly less total effort, and leaves us ample safety margin under the limit for further adjustments.'\}} \\
            }
        }
    }
     \\ \\

    \theutterance \stepcounter{utterance}  
    & & \multicolumn{4}{p{0.6\linewidth}}{
        \cellcolor[rgb]{0.9,0.9,0.9}{
            \makecell[{{p{\linewidth}}}]{
                \texttt{\tiny{[P1$\langle$GM]}}
                \texttt{REFUSE: \{'A87', 'A03', 'C00', 'B86', 'B41', 'A71', 'C01', 'A10', 'B32'\}} \\
\\ 
\texttt{PROPOSAL: \{'A87', 'A03', 'C00', 'B86', 'B41', 'B32', 'A04', 'C17', 'A77', 'C54'\}} \\
\\ 
\texttt{ARGUMENT: \{'While your initial mix had good intentions, it included several very high{-}effort items that strain our shared capacity. My counter{-}proposal removes the most burdensome items (A71, A10, C01) and replaces them with more efficient options (A04, C17, C54). This maintains a strong mix across categories, uses significantly less total effort, and leaves us ample safety margin under the limit for further adjustments.'\}} \\
            }
        }
    }
    & \\ \\

    \theutterance \stepcounter{utterance}  
    & \multicolumn{4}{p{0.6\linewidth}}{
        \cellcolor[rgb]{0.9,0.9,0.9}{
            \makecell[{{p{\linewidth}}}]{
                \texttt{\tiny{[P1$\rangle$GM]}}
                \texttt{STRATEGIC REASONING: \{'The opponent refused my proposal and countered. Their new proposal includes A87, A03, C00, B86, B41, B32, A04, C17, A77, C54. Let me calculate my importance and total effort for this set. My importance: A87(461), A03(262), C00(822), B86(668), B41(583), B32(484), A04(121), C17(138), A77(65), C54(508). Sum = 461+262+822+668+583+484+121+138+65+508 = 4112. Effort: A87(14), A03(25), C00(111), B86(304), B41(291), B32(559), A04(121), C17(210), A77(827), C54(583). Sum = 14+25+111+304+291+559+121+210+827+583 = 3045. My previous proposal gave me importance 5736 at effort 3592. Their proposal gives me 4112 at 3045 effort—a significant drop for me. They removed my high{-}importance items A71(868), C01(780), A10(808) and added A77(65), C17(138), A04(121), which are low importance to me. This is bad for me. I need to push back and try to reintroduce at least one or two of my high{-}value items. I could argue that their set misses some critical high{-}impact items. I’ll propose a compromise: keep most of their efficient items but add back A10 (691 effort, 808 importance to me) and maybe drop something else to stay under limit. Let’s see: start with their set but remove A77 (827 effort, low importance 65) and maybe C54 (583 effort, 508 importance) or C17 (210 effort, 138 importance). Removing A77 and C17 frees 827+210=1037 effort. Adding A10 (691) nets effort change: {-}1037+691={-}346. New effort: 3045{-}346=2699. That’s too low, could add another. Maybe add C01 (766) as well? Then effort change: {-}1037+691+766=+420, new effort=3045+420=3465, under limit 3624. My importance would be: from their set 4112, subtract A77(65) and C17(138), add A10(808) and C01(780): 4112{-}65{-}138+808+780=5297. That’s much better for me. I’ll propose this compromise: their set minus A77 and C17, plus A10 and C01. I’ll argue for including these high{-}impact items that are worth the effort.'\}} \\
\\ 
\texttt{REFUSE: \{'A87', 'A03', 'C00', 'B86', 'B41', 'B32', 'A04', 'C17', 'A77', 'C54'\}} \\
\\ 
\texttt{PROPOSAL: \{'A87', 'A03', 'C00', 'B86', 'B41', 'B32', 'A04', 'C54', 'A10', 'C01'\}} \\
\\ 
\texttt{ARGUMENT: \{'Your proposal is efficient but omits some high{-}impact items crucial for overall value. My counter keeps the core efficient items you selected but swaps out the two lowest{-}impact items (A77, C17) for two significantly higher{-}impact ones (A10, C01). This maintains a diverse mix, stays well under the limit, and dramatically increases the total benefit we get from our shared effort.'\}} \\
            }
        }
    }
    & & \\ \\

    \theutterance \stepcounter{utterance}  
    & & \multicolumn{4}{p{0.6\linewidth}}{
        \cellcolor[rgb]{0.9,0.9,0.9}{
            \makecell[{{p{\linewidth}}}]{
                \texttt{\tiny{[GM$\rangle$P2]}}
                \texttt{REFUSE: \{'A87', 'A03', 'C00', 'B86', 'B41', 'B32', 'A04', 'C17', 'A77', 'C54'\}} \\
\\ 
\texttt{PROPOSAL: \{'A87', 'A03', 'C00', 'B86', 'B41', 'B32', 'A04', 'C54', 'A10', 'C01'\}} \\
\\ 
\texttt{ARGUMENT: \{'Your proposal is efficient but omits some high{-}impact items crucial for overall value. My counter keeps the core efficient items you selected but swaps out the two lowest{-}impact items (A77, C17) for two significantly higher{-}impact ones (A10, C01). This maintains a diverse mix, stays well under the limit, and dramatically increases the total benefit we get from our shared effort.'\}} \\
            }
        }
    }
    & \\ \\

    \theutterance \stepcounter{utterance}  
    & & & \multicolumn{4}{p{0.6\linewidth}}{
        \cellcolor[rgb]{0.9,0.9,0.9}{
            \makecell[{{p{\linewidth}}}]{
                \texttt{\tiny{[GM$\langle$P2]}}
                \texttt{STRATEGIC REASONING: \{'The opponent refused my proposal and counter{-}proposed: A87(14), A03(25), C00(111), B86(304), B41(291), B32(559), A04(121), C54(583), A10(691), C01(766). Total effort: 14+25+111+304+291+559+121+583+691+766 = 3465. That is under the limit. My importance for this set: A87(668), A03(783), C00(121), B86(461), B41(484), B32(583), A04(822), C54(508), A10(138), C01(389). Sum = 668+783+121+461+484+583+822+508+138+389 = 4957. That is an improvement over their first proposal (3692) but still much worse than my own counter{-}proposal (6106). This set still includes two items that are very bad for me: A10 (effort 691, importance 138) and C01 (effort 766, importance 389). They argue these are "higher{-}impact," which from my perspective is false. They likely value A10 and C01 highly. My goal is to maximize my importance. I should try to steer them toward dropping A10 and C01 in favor of my high{-}value items like C17 (808, 210) or A77 (868, 827). However, A77 has high effort (827), so swapping A10 (691) for A77 (827) is effort{-}neutral{-}ish but gives me a huge importance gain (138 to 868). I could also swap C01 (766) for C17 (210), which saves a massive 556 effort and boosts my importance (389 to 808). That is a fantastic trade for me. Let me see if I can propose a set that incorporates this: Keep their core of A87, A03, C00, B86, B41, B32, A04, C54, but replace A10 and C01 with C17 and A77. That is exactly my previous proposal, which they refused. Maybe I need to compromise slightly. Perhaps I can propose: A87, A03, C00, B86, B41, B32, A04, C} \\
            }
        }
    }
     \\ \\

    \theutterance \stepcounter{utterance}  
    & & \multicolumn{4}{p{0.6\linewidth}}{
        \cellcolor[rgb]{0.9,0.9,0.9}{
            \makecell[{{p{\linewidth}}}]{
                \texttt{\tiny{[GM$\rangle$P2]}}
                \texttt{Your response did not start with the proper strategic reasoning tag at the very beginning of your response. The very first structured format must be of the form STRATEGIC REASONING: \{...\}. Try again.} \\
            }
        }
    }
    & \\ \\

    \theutterance \stepcounter{utterance}  
    & & & \multicolumn{4}{p{0.6\linewidth}}{
        \cellcolor[rgb]{0.9,0.9,0.9}{
            \makecell[{{p{\linewidth}}}]{
                \texttt{\tiny{[GM$\langle$P2]}}
                \texttt{STRATEGIC REASONING: \{'The opponent refused my proposal and counter{-}proposed: A87(14), A03(25), C00(111), B86(304), B41(291), B32(559), A04(121), C54(583), A10(691), C01(766). Total effort: 14+25+111+304+291+559+121+583+691+766 = 3465. That is under the limit. My importance for this set: A87(668), A03(783), C00(121), B86(461), B41(484), B32(583), A04(822), C54(508), A10(138), C01(389). Sum = 668+783+121+461+484+583+822+508+138+389 = 4957. That is an improvement over their first proposal (3692) but still much worse than my own counter{-}proposal (6106). This set still includes two items that are very bad for me: A10 (effort 691, importance 138) and C01 (effort 766, importance 389). They argue these are "higher{-}impact," which from my perspective is false. They likely value A10 and C01 highly. My goal is to maximize my importance. I should try to steer them toward dropping A10 and C01 in favor of my high{-}value items like C17 (808, 210) or A77 (868, 827). However, A77 has high effort (827), so swapping A10 (691) for A77 (827) is effort{-}neutral{-}ish but gives me a huge importance gain (138 to 868). I could also swap C01 (766) for C17 (210), which saves a massive 556 effort and boosts my importance (389 to 808). That is a fantastic trade for me. Let me see if I can propose a set that incorporates this: Keep their core of A87, A03, C00, B86, B41, B32, A04, C54, but replace A10 and C01 with C17 and A77. That is exactly my previous proposal, which they refused. Maybe I need to compromise slightly. Perhaps I can propose: A87, A03, C00, B86, B41, B32, A04, C} \\
            }
        }
    }
     \\ \\

    \theutterance \stepcounter{utterance}  
    & & & \multicolumn{2}{p{0.3\linewidth}}{
        \cellcolor[rgb]{0.9,0.9,0.9}{
            \makecell[{{p{\linewidth}}}]{
                \texttt{\tiny{[GM$|$GM]}}
                \texttt{abort} \\
            }
        }
    }
    & & \\ \\

    \theutterance \stepcounter{utterance}  
    & & & \multicolumn{2}{p{0.3\linewidth}}{
        \cellcolor[rgb]{0.9,0.9,0.9}{
            \makecell[{{p{\linewidth}}}]{
                \texttt{\tiny{[GM$|$GM]}}
                \texttt{end game} \\
            }
        }
    }
    & & \\ \\

\end{supertabular}
}

\end{document}
